%% 由 build/emit_tex.py 生成（生成物，勿手改）；定制请放 user_style.tex / user_meta.yaml
\chapter{全书总参考文献（附录
G）}\label{ux5168ux4e66ux603bux53c2ux8003ux6587ux732eux9644ux5f55-g}

\begin{quote}
本页由 \texttt{build/gen\_references.py}
\textbf{自动生成}，勿手改。单一事实源 = 各章 \texttt{references.md} +
附录 \texttt{references.md}；新增/删除引用后重新运行生成即可。
速查方式：按类别浏览；每行``用于''列标注了哪些章节使用该资料，点章节链接可回正文。
\end{quote}

共汇总 \textbf{163 条}（来自 9 个章节 + 附录）。

\section{一、官方文档与工具链}\label{ux4e00ux5b98ux65b9ux6587ux6863ux4e0eux5de5ux5177ux94fe}

{\def\LTcaptype{none} % do not increment counter
\begin{longtable}[]{@{}
  >{\raggedright\arraybackslash}p{(\linewidth - 4\tabcolsep) * \real{0.3333}}
  >{\raggedright\arraybackslash}p{(\linewidth - 4\tabcolsep) * \real{0.3333}}
  >{\raggedright\arraybackslash}p{(\linewidth - 4\tabcolsep) * \real{0.3333}}@{}}
\toprule\noalign{}
\begin{minipage}[b]{\linewidth}\raggedright
资料
\end{minipage} & \begin{minipage}[b]{\linewidth}\raggedright
说明
\end{minipage} & \begin{minipage}[b]{\linewidth}\raggedright
用于
\end{minipage} \\
\midrule\noalign{}
\endhead
\bottomrule\noalign{}
\endlastfoot
\href{https://www.statsmodels.org/stable/api.html}{API Reference} &
全部函数索引 & 第7章 Statsmodels \\
\href{https://www.statsmodels.org/stable/generated/statsmodels.tsa.arima.model.ARIMA.html}{ARIMA}
& ARIMA/SARIMAX 模型 & 第7章 Statsmodels \\
\href{https://www.anaconda.com/docs/getting-started/concepts/anaconda-or-miniconda}{Anaconda
vs Miniconda} & 选型说明（课程推荐 Miniconda） & 第0章 前置基础 \\
\href{https://numpy.org/doc/stable/user/basics.broadcasting.html}{Broadcasting
官方说明} & 广播规则权威定义 & 第1章 NumPy、附录 \\
\href{https://networkx.org/documentation/stable/auto_examples/index.html}{Complex
Network Analysis in Python（官方示例）} & 官方示例集，覆盖绘图/社区/路径
& 第6章 NetworkX \\
\href{https://www.statsmodels.org/stable/examples/index.html}{Examples
总目录} & 官方示例（含回归与时间序列） & 第7章 Statsmodels \\
\href{https://www.statsmodels.org/stable/formula.html}{Formula API} &
ols(`y \textasciitilde{} x', data=\ldots) 语法 & 第7章 Statsmodels \\
\href{https://www.statsmodels.org/stable/generated/statsmodels.genmod.generalized_linear_model.GLM.html}{GLM}
& 广义线性模型 & 第7章 Statsmodels \\
\href{https://git-scm.com/doc}{Git 官方文档} & 书/命令/教程 & 第0章
前置基础 \\
\href{https://numpy.org/doc/stable/reference/routines.io.html}{I/O 参考}
& save/loadtxt/savez & 第1章 NumPy \\
\href{https://docs.jupyter.org/en/latest/}{Jupyter 官方文档} &
笔记本/内核/导出 & 第0章 前置基础 \\
\href{https://www.learnlatex.org/zh-hans/}{Learn LaTeX 中文站} &
免费结构化课程：结构/公式/表格/插图 & 第0章 前置基础 \\
\href{https://www.loongtex.com/}{LoongTeX（龙文）} &
国产在线编译（备选，中文友好） & 第0章 前置基础 \\
\href{https://matplotlib.org/stable/cheatsheets/index.html}{Matplotlib
Cheatsheets} & 速查表（含多子图排版） & 第5章 Matplotlib \\
\href{https://matplotlib.org/stable/gallery/index.html}{Matplotlib
图型画廊} & 各类图型官方示例 & 第5章 Matplotlib \\
\href{https://matplotlib.org/stable/tutorials/introductory/quick_start.html}{Matplotlib
快速开始} & 官方绝对入门，覆盖 Plot/Axes/Figure & 第5章 Matplotlib \\
\href{https://matplotlib.org/stable/users/explain/text/text_props.html}{Matplotlib
文本/字体} & 字体、字号、数学公式 & 第5章 Matplotlib \\
\href{https://matplotlib.org/stable/tutorials/intermediate/tight_layout_guide.html}{Matplotlib:
Anatomy of a Figure} & 子图排版与 tight\_layout 细节 & 第5章
Matplotlib \\
\href{https://networkx.org/documentation/stable/release/release_3.3.html}{NetworkX
3.x 迁移/新特性} & 版本更新说明 & 第6章 NetworkX \\
\href{https://networkx.org/documentation/stable/tutorial.html}{NetworkX
Tutorial（官方教程）} & 从零创建图、访问节点边、常用操作 & 第6章
NetworkX、附录 \\
\href{https://networkx.cn/documentation/stable/}{NetworkX 中文文档镜像}
& 中文版官方文档（有翻译） & 第6章 NetworkX \\
\href{https://networkx.org/documentation/stable/faq.html}{NetworkX
社区（Discussions/Stack Overflow）} & 常见问题与讨论入口 & 第6章
NetworkX \\
\href{https://numpy.org/doc/stable/user/quickstart.html}{NumPy
Quickstart} & 快速上手教程 & 第1章 NumPy \\
\href{https://numpy.org/doc/stable/user/numpy-for-matlab-users.html}{NumPy
for MATLAB users} & 给 MATLAB 背景学生迁移 & 第1章 NumPy \\
\href{https://numpy.org/doc/stable/user/absolute_beginners.html}{NumPy：Absolute
Beginner} & 官方``绝对入门''，覆盖创建/索引/广播 & 第1章 NumPy \\
\href{https://www.statsmodels.org/stable/generated/statsmodels.regression.linear_model.OLS.html}{OLS}
& 最小二乘回归 & 第7章 Statsmodels \\
\href{https://www.overleaf.com/}{Overleaf} & 在线编译（备选） & 第0章
前置基础 \\
\href{https://pandas.pydata.org/docs/user_guide/10min.html}{Pandas 10
分钟入门} & 一分钟建立 Series/DataFrame 心智 & 第4章 Pandas \\
\href{https://pandas.pydata.org/Pandas_Cheat_Sheet.pdf}{Pandas Cheat
Sheet（官方）} & 一页速查 & 第4章 Pandas \\
\href{https://pandas.pydata.org/docs/}{Pandas
中文入门（pandas.pydata.org 官方中文翻译入口）} & 官方文档，含中文化镜像
& 第4章 Pandas \\
\href{https://pandas.pydata.org/}{Pandas 官方主页} & 库主页、版本、安装
& 第4章 Pandas \\
\href{https://pandas.pydata.org/docs/user_guide/index.html}{Pandas
用户指南} & 索引/清洗/分组/时间序列全指南 & 第4章 Pandas \\
\href{https://git-scm.com/book/zh/v2}{Pro Git 中文版} & 系统化学习 &
第0章 前置基础 \\
\href{https://matplotlib.org/stable/tutorials/pyplot.html}{Pyplot 教程}
& 函数式接口与状态机说明 & 第5章 Matplotlib \\
\href{https://docs.python.org/zh-cn/3/tutorial/}{Python
官方教程（中文）} & 语法查漏的第一权威 & 第0章 前置基础 \\
\href{https://www.statsmodels.org/stable/regression.html}{Regressions} &
线性回归/GLS/WLS 总览 & 第7章 Statsmodels \\
\href{https://docs.scipy.org.cn/doc/scipy/}{SciPy 中文文档镜像} &
中文版参考（版本可能较旧） & 第3章 SciPy \\
\href{https://scipy.org/}{SciPy 主页} & 库介绍、版本、安装 & 第3章
SciPy \\
\href{https://docs.scipy.org/doc/scipy/reference/}{SciPy
参考手册（全模块）} & 按模块查函数签名与示例 & 第3章 SciPy、附录 \\
\href{https://docs.scipy.org/doc/scipy/tutorial/}{SciPy
官方教程（tutorial 目录）} & 官方讲解积分/FFT/插值 & 第3章 SciPy \\
\href{https://docs.scipy.org/doc/scipy/release.html}{SciPy 版本发布说明}
& 查看函数改名/废弃 & 第3章 SciPy \\
\href{https://seaborn.pydata.org/tutorial.html}{Seaborn 官方教程} &
分布/箱线/热力图/分面教程 & 第5章 Matplotlib \\
\href{https://github.com/ljtlrh/statsmodels_doc_zh}{Statsmodels
中文文档（ApacheCN 翻译）} & 官方文档的中文翻译 & 第7章 Statsmodels \\
\href{https://www.statsmodels.org/stable/index.html}{Statsmodels 主页} &
库简介与总入口 & 第7章 Statsmodels、附录 \\
\href{https://www.statsmodels.org/stable/examples/index.html\#regression}{Statsmodels
官方示例（回归）} & 一整套可运行回归案例 & 第7章 Statsmodels \\
\href{https://www.statsmodels.org/stable/examples/index.html\#time-series}{Statsmodels
官方示例（时间序列）} & 分解、ARIMA/SARIMA 案例 & 第7章 Statsmodels \\
\href{https://docs.sympy.org/latest/tutorials/intro-tutorial/calculus.html}{SymPy
Calculus（求导/积分/极限）} & 极限、diff、integrate 的权威说明 & 第2章
SymPy \\
\href{https://docs.sympy.org/latest/tutorial/gotchas.html}{SymPy
Gotchas（常见坑）} & 官方 FAQ，避开符号计算陷阱 & 第2章 SymPy \\
\href{https://docs.sympy.org/latest/modules/matrices/index.html}{SymPy
Matrices} & 符号矩阵、特征值、分解 & 第2章 SymPy \\
\href{https://docs.sympy.org/latest/modules/solvers/ode.html}{SymPy ODE
(dsolve)} & 微分方程符号求解 & 第2章 SymPy \\
\href{https://docs.sympy.org/latest/modules/solvers/index.html}{SymPy
Solvers} & solve/solveset/linsolve/nonlinsolve 参考 & 第2章 SymPy \\
\href{https://docs.sympy.org/latest/tutorials/index.html}{SymPy
Tutorials（总目录）} & 官方入门教程入口，含基础与进阶 & 第2章 SymPy \\
\href{https://docs.sympy.org/latest/modules/utilities/lambdify.html}{SymPy
lambdify} & 符号→数值函数的关键工具 & 第2章 SymPy \\
\href{https://docs.sympy.org/latest/index.html}{SymPy 官方 1.13 文档} &
完整 API 参考 & 第2章 SymPy \\
\href{https://docs.sympy.org/latest/tutorials/intro-tutorial/}{SymPy
官方 doc 中的 Examples} & 每条均可上机复现 & 第2章 SymPy \\
\href{https://www.statsmodels.org/stable/tsa.html}{TSA（时间序列）} &
时间序列模块总览 & 第7章 Statsmodels \\
\href{https://www.tug.org/texlive/acquire-netinstall.html}{TeX Live
下载} & Windows 安装器 & 第0章 前置基础 \\
\href{https://www.statsmodels.org/stable/user-guide.html}{User Guide} &
全书式教学指南 & 第7章 Statsmodels \\
\href{https://code.visualstudio.com/docs/python/python-tutorial}{VS Code
Python 入门} & 从新建文件到运行调试 & 第0章 前置基础 \\
\href{https://code.visualstudio.com/docs/python/environments}{VS Code
Python 环境官方文档} & 解释器/虚拟环境/内核原理 & 第0章 前置基础 \\
\href{https://code.visualstudio.com/}{VS Code 官方下载} & 编辑器 & 第0章
前置基础 \\
\href{https://www.statsmodels.org/stable/generated/statsmodels.regression.linear_model.WLS.html}{WLS}
& 加权最小二乘 & 第7章 Statsmodels \\
\href{https://www.statsmodels.org/stable/generated/statsmodels.tsa.stattools.adfuller.html}{adfuller}
& ADF 单位根检验 & 第7章 Statsmodels \\
\href{https://www.statsmodels.org/stable/generated/statsmodels.stats.anova.anova_lm.html}{anova\_lm}
& 方差分析 & 第7章 Statsmodels \\
\href{https://ctan.org/pkg/ctex}{ctex 宏包文档} & 中文排版权威 & 第0章
前置基础 \\
\href{https://numpy.org/doc/stable/user/basics.html}{ndarray 基础} &
形状、索引、切片、视图/拷贝 & 第1章 NumPy \\
\href{https://numpy.org/doc/stable/reference/routines.linalg.html}{numpy.linalg
参考} & 线性代数全部函数 & 第1章 NumPy、附录 \\
\href{https://numpy.org/doc/stable/reference/routines.polynomials.html}{numpy.polynomial
参考} & 多项式接口 & 第1章 NumPy \\
\href{https://pandas.pydata.org/docs/reference/api/pandas.DataFrame.html}{pandas.DataFrame
参考} & DataFrame 全部方法 & 第4章 Pandas \\
\href{https://pandas.pydata.org/docs/reference/api/pandas.Series.html}{pandas.Series
参考} & Series 全部方法 & 第4章 Pandas \\
\href{https://pandas.pydata.org/docs/reference/api/pandas.pivot_table.html}{pivot\_table
参考} & 透视表 & 第4章 Pandas \\
\href{https://scikit-learn.org/stable/modules/classes.html}{scikit-learn
API 速查} & 全部类/函数索引 & 第8章 scikit-learn \\
\href{https://inria.github.io/scikit-learn-mooc/}{scikit-learn
MOOC（INRIA）} & 官方协作的免费课程，带习题与 Notebook & 第8章
scikit-learn \\
\href{https://scikit-learn.org/stable/modules/cross_validation.html}{scikit-learn
交叉验证} & train/test、交叉验证、调参 & 第8章 scikit-learn \\
\href{https://scikit-learn.org/stable/tutorial/index.html}{scikit-learn
入门教程} & 官方快速上手，含``选择正确评估器''参考 & 第8章
scikit-learn \\
\href{https://scikit-learn.org/stable/}{scikit-learn 官方主页 /
稳定版文档} & 权威 API 与用户指南入口，建议作为查错第一站 & 第8章
scikit-learn \\
\href{https://scikit-learn.org/stable/auto_examples/index.html}{scikit-learn
官方示例库} & 大量可运行示例，按算法/任务索引 & 第8章 scikit-learn \\
\href{https://scikit-learn.org/stable/modules/model_evaluation.html}{scikit-learn
模型评估} & 分类/回归/聚类指标 & 第8章 scikit-learn \\
\href{https://scikit-learn.org/stable/modules/feature_selection.html}{scikit-learn
特征选择} & 过滤/包装/嵌入三类方法 & 第8章 scikit-learn \\
\href{https://scikit-learn.org/stable/user_guide.html}{scikit-learn
用户指南} & 按主题讲解预处理、监督、无监督、模型选择 & 第8章
scikit-learn、附录 \\
\href{https://scikit-learn.org/stable/modules/preprocessing.html}{scikit-learn
预处理} & 编码、缩放、缺失值处理 & 第8章 scikit-learn \\
\href{https://docs.scipy.org/doc/scipy/reference/fft.html}{scipy.fft
傅里叶变换} & \texttt{fft}/\texttt{rfft}/\texttt{fftfreq} & 第3章
SciPy \\
\href{https://docs.scipy.org/doc/scipy/reference/integrate.html}{scipy.integrate
积分与 ODE} &
\texttt{quad}/\texttt{odeint}/\texttt{solve\_ivp}/\texttt{solve\_bvp} &
第3章 SciPy \\
\href{https://docs.scipy.org/doc/scipy/reference/interpolate.html}{scipy.interpolate
插值} & 一维/多维/散点插值 & 第3章 SciPy \\
\href{https://docs.scipy.org/doc/scipy/reference/optimize.html}{scipy.optimize
优化} & 求根/极值/规划/拟合 & 第3章 SciPy \\
\href{https://docs.scipy.org/doc/scipy/reference/signal.html}{scipy.signal
信号处理} & 滤波器设计、卷积 & 第3章 SciPy \\
\href{https://docs.scipy.org/doc/scipy/reference/stats.html}{scipy.stats
统计} & 检验/分布/描述统计 & 第3章 SciPy \\
\href{https://www.statsmodels.org/stable/generated/statsmodels.tsa.seasonal.seasonal_decompose.html}{seasonal\_decompose}
& 时间序列分解 & 第7章 Statsmodels \\
\href{https://www.statsmodels.org/stable/}{statsmodels 官方} &
方差分析事后比较/回归 & 第3章 SciPy \\
\href{https://networkx.org/documentation/stable/reference/algorithms/centrality.html}{中心性算法}
& 介数/接近/度中心性 & 第6章 NetworkX \\
\href{https://pandas.pydata.org/docs/user_guide/groupby.html}{分组（groupby）}
& 分组聚合/变换/过滤 & 第4章 Pandas \\
\href{https://networkx.org/documentation/stable/reference/convert.html}{图与
NumPy/SciPy 互转} & to\_numpy\_array / to\_scipy\_sparse\_array & 第6章
NetworkX \\
\href{https://networkx.org/documentation/stable/reference/generators.html}{图生成器}
& path/cycle/complete/karate/随机图 & 第6章 NetworkX \\
\href{https://networkx.org/documentation/stable/reference/classes/index.html}{图类型（Graph
/ DiGraph 等）} & Graph / DiGraph / MultiGraph 接口 & 第6章 NetworkX \\
\href{https://docs.conda.io/en/latest/miniconda.html}{安装 Miniconda} &
官方下载入口（速度慢用清华镜像） & 第0章 前置基础 \\
\href{https://networkx.org/documentation/stable/auto_examples/graph/plot_karate_club.html}{官方示例：空手道俱乐部}
& 经典社区网络可视化 & 第6章 NetworkX \\
\href{https://pandas.pydata.org/docs/user_guide/enhancingperf.html}{性能提升}
& 向量化、chunksize 等 & 第4章 Pandas \\
\href{https://www.statsmodels.org/devel/examples/notebooks/generated/stl_decomposition.html}{时间序列
STL 示例} & LOESS 季节性-趋势分解 & 第7章 Statsmodels \\
\href{https://pandas.pydata.org/docs/user_guide/timeseries.html}{时间序列指南}
& date\_range/resample/rolling & 第4章 Pandas \\
\href{https://networkx.org/documentation/stable/reference/algorithms/tree.html}{最小生成树}
& minimum\_spanning\_tree 等 & 第6章 NetworkX \\
\href{https://networkx.org/documentation/stable/reference/algorithms/shortest_paths.html}{最短路径算法}
& Dijkstra / Bellman-Ford / Floyd-Warshall & 第6章 NetworkX \\
\href{https://networkx.org/documentation/stable/reference/algorithms/community.html}{社区检测与模块度}
& greedy\_modularity / louvain / modularity & 第6章 NetworkX \\
\href{https://networkx.org/documentation/stable/reference/algorithms/index.html}{算法总览}
& 全部算法模块入口 & 第6章 NetworkX \\
\href{https://networkx.org/documentation/stable/reference/drawing.html}{绘图
API} & draw / spring\_layout / 自定义 & 第6章 NetworkX \\
\href{https://pandas.pydata.org/docs/user_guide/missing_data.html}{缺失值处理}
& isna/fillna/dropna/ffill/bfill & 第4章 Pandas \\
\href{https://networkx.org/documentation/stable/reference/algorithms/connectivity.html}{连通性}
& is\_connected / 强连通 / 分量 & 第6章 NetworkX \\
\href{https://numpy.org/doc/stable/reference/random/generator.html}{随机数（Generator）}
& \texttt{default\_rng} 用法 & 第1章 NumPy \\
\href{https://matplotlib.org/stable/tutorials/introductory/usage.html\#the-object-oriented-interface}{面向对象接口}
& Figure/Axes 对象模型权威定义 & 第5章 Matplotlib \\
\end{longtable}
}

\section{二、教材与书籍}\label{ux4e8cux6559ux6750ux4e0eux4e66ux7c4d}

{\def\LTcaptype{none} % do not increment counter
\begin{longtable}[]{@{}
  >{\raggedright\arraybackslash}p{(\linewidth - 4\tabcolsep) * \real{0.3333}}
  >{\raggedright\arraybackslash}p{(\linewidth - 4\tabcolsep) * \real{0.3333}}
  >{\raggedright\arraybackslash}p{(\linewidth - 4\tabcolsep) * \real{0.3333}}@{}}
\toprule\noalign{}
\begin{minipage}[b]{\linewidth}\raggedright
资料
\end{minipage} & \begin{minipage}[b]{\linewidth}\raggedright
说明
\end{minipage} & \begin{minipage}[b]{\linewidth}\raggedright
用于
\end{minipage} \\
\midrule\noalign{}
\endhead
\bottomrule\noalign{}
\endlastfoot
\href{https://space.bilibili.com/88461692/}{3Blue1Brown
线性代数的本质（中文字幕）} & 可视化直觉 & 附录 \\
\href{https://github.com/datawhalechina/grape-book}{Datawhale
图深度学习（葡萄书）} & 中文开源书，进阶图神经网络 & 第6章 NetworkX \\
\href{https://mml-book.github.io/}{Mathematics for Machine Learning} &
Part I 数学基础（线性代数/微积分）+ Part II 机器学习，主题式 & 附录 \\
\href{https://jakevdp.github.io/PythonDataScienceHandbook/}{Python Data
Science Handbook} & Jake VanderPlas 的《线性回归》一章对比
sklearn/statsmodels & 第7章 Statsmodels \\
\href{https://github.com/jakevdp/PythonDataScienceHandbook}{Python Data
Science Handbook (Jake VanderPlas)} & 第 2 章 NumPy，可整章对照 & 第1章
NumPy、第2章 SymPy、第3章 SciPy、第4章 Pandas、第5章 Matplotlib、第6章
NetworkX、第8章 scikit-learn \\
\href{https://github.com/AllenDowney/ThinkStats2}{thinkstats2} & 用
Python 讲统计 & 附录 \\
\href{https://github.com/datawhalechina/pumpkin-book}{南瓜书（《机器学习》公式推导）}
& 与周志华《机器学习》配套的公式推导详解 & 第8章 scikit-learn \\
\href{https://github.com/datawhalechina/intro-mathmodel}{数学建模导论} &
后续课程 & 第1章 NumPy、第3章 SciPy、第4章 Pandas、附录 \\
\href{https://github.com/datawhalechina/scientific-computing}{本项目（科学计算）}
& 本书开源仓库 & 第1章 NumPy、第2章 SymPy、第3章 SciPy、第4章
Pandas、第5章 Matplotlib、第7章 Statsmodels、第8章 scikit-learn \\
\href{https://github.com/datawhalechina/leedl-tutorial}{李宏毅机器学习（Datawhale
整理）} & 中文讲义，串起 ML 与 DL 的``为什么'' & 第8章 scikit-learn \\
\end{longtable}
}

\section{三、精品教程与课程}\label{ux4e09ux7cbeux54c1ux6559ux7a0bux4e0eux8bfeux7a0b}

{\def\LTcaptype{none} % do not increment counter
\begin{longtable}[]{@{}
  >{\raggedright\arraybackslash}p{(\linewidth - 4\tabcolsep) * \real{0.3333}}
  >{\raggedright\arraybackslash}p{(\linewidth - 4\tabcolsep) * \real{0.3333}}
  >{\raggedright\arraybackslash}p{(\linewidth - 4\tabcolsep) * \real{0.3333}}@{}}
\toprule\noalign{}
\begin{minipage}[b]{\linewidth}\raggedright
资料
\end{minipage} & \begin{minipage}[b]{\linewidth}\raggedright
说明
\end{minipage} & \begin{minipage}[b]{\linewidth}\raggedright
用于
\end{minipage} \\
\midrule\noalign{}
\endhead
\bottomrule\noalign{}
\endlastfoot
\href{https://github.com/rougier/numpy-100}{100 NumPy
Exercises（rougier）} & 100 题带解答，适合课后 & 第1章 NumPy \\
\href{https://github.com/datawhalechina/fantastic-matplotlib}{Datawhale
fantastic-matplotlib} & 中文、渐进式 Matplotlib 教程 & 第5章
Matplotlib \\
\href{https://github.com/datawhalechina/wow-plotly}{Datawhale
wow-plotly} & Plotly 交互式教程（对比阅读） & 第5章 Matplotlib \\
\href{https://blog.csdn.net/weixin_45306755/article/details/108122794}{Datawhale
时间序列相关笔记} & 时间序列特征与规则入门 & 第7章 Statsmodels \\
\href{https://github.com/ageron/handson-ml3}{Hands-On ML 3rd（Aurélien
Géron）} & 第 2 章起系统讲 sklearn，示例丰富 & 第8章 scikit-learn \\
\href{https://www.kaggle.com/learn/intro-to-machine-learning}{Kaggle
Intro to ML} & 用真实数据练决策树/随机森林/交叉验证 & 第8章
scikit-learn \\
\href{https://learnxinyminutes.com/docs/zh-cn/python3-cn/}{Learn X in Y
minutes: Python} & 半小时过一遍语法 & 第0章 前置基础 \\
\href{https://github.com/jrjohansson/scientific-python-lectures}{Lectures
on Scientific Computing (Robert Johansson)} &
Lecture-2-Numpy.ipynb，Notebook 形式 & 第1章 NumPy、第2章 SymPy、第3章
SciPy、第4章 Pandas、第6章 NetworkX、第7章 Statsmodels \\
\href{https://ocw.mit.edu/courses/2-086-numerical-computation-for-mechanical-engineers-spring-2013/}{MIT
OCW 数值计算（2.086）} & 数值方法讲义 & 附录 \\
\href{https://numpy.net.cn/learn/}{NumPy 中文学习站} &
中文入门与文档镜像 & 第1章 NumPy \\
\href{https://liaoxuefeng.com/books/python/introduction/}{Python 教程 ·
廖雪峰（中文）} & 中文查漏 & 第0章 前置基础 \\
\href{https://www.zhihu.com/question/582168949}{Python数据可视化（知乎/博客）}
& 中文快速入门与经验 & 第5章 Matplotlib \\
\href{https://scipy-lectures.org/}{SciPy Lecture Notes} & 系统、带图，含
SciPy 章节 & 第3章 SciPy、第4章 Pandas、第6章 NetworkX、附录 \\
\href{https://scipy-lectures.org/packages/10_advanced_2d_plotting.html}{SciPy
Lecture Notes} & Matplotlib 进阶绘图讲义 & 第5章 Matplotlib \\
\href{https://scipy-lectures.org/intro/numpy/index.html}{SciPy Lecture
Notes -- NumPy} & 系统、带图，覆盖数组/广播/运算 & 第1章 NumPy \\
\href{https://scipy-lectures.org/packages/sympy.html}{SciPy Lectures --
SymPy} & 系统、带例子，覆盖符号计算核心 & 第2章 SymPy \\
\href{https://zetcode.cn/python/sympy/}{ZetCode 中文教程：SymPy} &
中文入门，含符号运算示例 & 第2章 SymPy \\
\href{https://numpy.net.cn/}{numpy.net.cn（中文学习站）} &
中文科学计算入门（含 Pandas 相关） & 第4章 Pandas \\
\href{https://github.com/casperdoudou/sklearn-doc-zh}{sklearn
中文文档（社区译本）} & 中文版用户指南/API 参考，适合初读 & 第8章
scikit-learn \\
\href{https://blog.csdn.net/Rocky006/article/details/148866312}{华为云/CSDN
NetworkX 教程} & 中文使用详解，含代码 & 第6章 NetworkX \\
\href{https://zh.khanacademy.org/}{可汗学院（中文）} & 微积分/统计基础 &
附录 \\
\href{https://liaoxuefeng.com/books/git/introduction/}{廖雪峰 Git
教程（中文）} & 中文入门 & 第0章 前置基础 \\
\href{https://cloud.tencent.cn/developer/article/2449258}{猫头虎 SciPy
入门教程} & 安装/模块/用例中文入门 & 第3章 SciPy \\
\href{https://zhuanlan.zhihu.com/p/111573239}{知乎专栏：SymPy
入门与实战} & 中文文章，带代码与截图 & 第2章 SymPy \\
\href{https://zhuanlan.zhihu.com/p/591617257}{知乎：复杂网络建模（Python+NetworkX）}
& 复杂网络建模课程代码（中文） & 第6章 NetworkX \\
\href{https://github.com/datawhalechina/learn-python-the-smart-way-v2}{聪明办法学
Python v2（Datawhale）} & 前置课程：Chap0 安装、Chap1 启航、Chap2-6
基础语法 & 第0章 前置基础、第1章 NumPy、第2章 SymPy、第3章 SciPy、第4章
Pandas、第5章 Matplotlib、第7章 Statsmodels、第8章 scikit-learn \\
\end{longtable}
}

\section{四、习题与实战}\label{ux56dbux4e60ux9898ux4e0eux5b9eux6218}

{\def\LTcaptype{none} % do not increment counter
\begin{longtable}[]{@{}
  >{\raggedright\arraybackslash}p{(\linewidth - 4\tabcolsep) * \real{0.3333}}
  >{\raggedright\arraybackslash}p{(\linewidth - 4\tabcolsep) * \real{0.3333}}
  >{\raggedright\arraybackslash}p{(\linewidth - 4\tabcolsep) * \real{0.3333}}@{}}
\toprule\noalign{}
\begin{minipage}[b]{\linewidth}\raggedright
资料
\end{minipage} & \begin{minipage}[b]{\linewidth}\raggedright
说明
\end{minipage} & \begin{minipage}[b]{\linewidth}\raggedright
用于
\end{minipage} \\
\midrule\noalign{}
\endhead
\bottomrule\noalign{}
\endlastfoot
\href{https://github.com/ajcr/100-pandas-puzzles}{100 Pandas Puzzles} &
100 道 Pandas 小题带答案 & 第4章 Pandas \\
\href{https://www.ntnu.no/wiki/spaces/imtsoftware/pages/148770976/Exercises+and+solutions+symbolic+mathematics+in+Python}{NTNU：Exercises
and solutions in symbolic mathematics (SymPy)} & 含 SymPy
符号计算练习题与答案 & 第2章 SymPy \\
\href{https://github.com/fengdu78/Data-Science-Notes}{NumPy
习题（中文版，Data-Science-Notes）} & 含 numpy-100 中文整理 & 第1章
NumPy \\
\href{https://github.com/guipsamora/pandas_exercises}{Pandas Exercises
(guipsamora)} & 分主题练习（清洗/分组/时间序列） & 第4章 Pandas \\
\href{https://scipy.github.io/old-wiki/pages/Cookbook/}{SciPy Cookbook}
& 大量实战小例子（老但仍有效） & 第3章 SciPy \\
\href{https://github.com/4GeeksAcademy/numpy-100}{numpy\_exercises by
4GeeksAcademy} & 在线运行版本（Binder 友好） & 第1章 NumPy \\
\end{longtable}
}

\section{五、开源项目与配套仓库}\label{ux4e94ux5f00ux6e90ux9879ux76eeux4e0eux914dux5957ux4ed3ux5e93}

{\def\LTcaptype{none} % do not increment counter
\begin{longtable}[]{@{}
  >{\raggedright\arraybackslash}p{(\linewidth - 4\tabcolsep) * \real{0.3333}}
  >{\raggedright\arraybackslash}p{(\linewidth - 4\tabcolsep) * \real{0.3333}}
  >{\raggedright\arraybackslash}p{(\linewidth - 4\tabcolsep) * \real{0.3333}}@{}}
\toprule\noalign{}
\begin{minipage}[b]{\linewidth}\raggedright
资料
\end{minipage} & \begin{minipage}[b]{\linewidth}\raggedright
说明
\end{minipage} & \begin{minipage}[b]{\linewidth}\raggedright
用于
\end{minipage} \\
\midrule\noalign{}
\endhead
\bottomrule\noalign{}
\endlastfoot
\href{https://github.com/wesm/pydata-book}{Python for Data Analysis 3rd
(Wes McKinney)} & Pandas 作者亲著，第三版配套代码 & 第4章 Pandas、第7章
Statsmodels \\
\href{https://github.com/dair-ai/Mathematics-for-ML}{dair-ai/Mathematics-for-ML}
& 按主题收集资源的清单 & 附录 \\
\href{https://github.com/tomdonald/matplotlib-cn}{matplotlib-cn（中文文档）}
& Matplotlib 中文文档镜像 & 第5章 Matplotlib \\
\href{https://github.com/mint-lab/prog_meets_math}{mint-lab/prog\_meets\_math}
& ``Python Meets Math''：Calculus / Linear Algebra / Optimization /
Probability + 代码 & 附录 \\
\end{longtable}
}

\section{六、中文资料与镜像}\label{ux516dux4e2dux6587ux8d44ux6599ux4e0eux955cux50cf}

{\def\LTcaptype{none} % do not increment counter
\begin{longtable}[]{@{}
  >{\raggedright\arraybackslash}p{(\linewidth - 4\tabcolsep) * \real{0.3333}}
  >{\raggedright\arraybackslash}p{(\linewidth - 4\tabcolsep) * \real{0.3333}}
  >{\raggedright\arraybackslash}p{(\linewidth - 4\tabcolsep) * \real{0.3333}}@{}}
\toprule\noalign{}
\begin{minipage}[b]{\linewidth}\raggedright
资料
\end{minipage} & \begin{minipage}[b]{\linewidth}\raggedright
说明
\end{minipage} & \begin{minipage}[b]{\linewidth}\raggedright
用于
\end{minipage} \\
\midrule\noalign{}
\endhead
\bottomrule\noalign{}
\endlastfoot
\href{https://datawhalechina.github.io/grape-book}{Grape-book（图深度学习）}
& 图神经网络进阶（延伸） & 第6章 NetworkX \\
\href{https://mirrors.tuna.tsinghua.edu.cn/}{清华 Anaconda/PyPI 镜像} &
conda/pip 都从这里加速 & 第0章 前置基础 \\
\href{https://mirrors.tuna.tsinghua.edu.cn/anaconda/miniconda/}{清华
Miniconda 镜像} & 国内下载加速 & 第0章 前置基础 \\
\href{https://cloud.tencent.cn/developer/article/2510103}{腾讯云开发者社区：Python
科学计算之 SymPy} & 中文案例集合 & 第2章 SymPy \\
\end{longtable}
}

\section{七、其他}\label{ux4e03ux5176ux4ed6}

{\def\LTcaptype{none} % do not increment counter
\begin{longtable}[]{@{}
  >{\raggedright\arraybackslash}p{(\linewidth - 4\tabcolsep) * \real{0.3333}}
  >{\raggedright\arraybackslash}p{(\linewidth - 4\tabcolsep) * \real{0.3333}}
  >{\raggedright\arraybackslash}p{(\linewidth - 4\tabcolsep) * \real{0.3333}}@{}}
\toprule\noalign{}
\begin{minipage}[b]{\linewidth}\raggedright
资料
\end{minipage} & \begin{minipage}[b]{\linewidth}\raggedright
说明
\end{minipage} & \begin{minipage}[b]{\linewidth}\raggedright
用于
\end{minipage} \\
\midrule\noalign{}
\endhead
\bottomrule\noalign{}
\endlastfoot
\href{https://www.bigocheatsheet.com/}{Big-O Cheat Sheet} &
常用复杂度速查 & 附录 \\
\href{https://scipy.github.io/old-wiki/pages/Cookbook/OptimizationAndFitDemo1.html}{Optimization
and Fit Cookbook} & 拟合示例详解 & 第3章 SciPy \\
\href{https://www.pythoncheatsheet.org/}{Python Cheat Sheet} & 速查表 &
第0章 前置基础 \\
\href{https://seaborn.pydata.org/api.html}{Seaborn API} & 函数索引 &
第5章 Matplotlib \\
\href{https://seaborn.pydata.org/examples/index.html}{Seaborn 示例} &
多分组/分面/热力图示例 & 第5章 Matplotlib \\
\end{longtable}
}

\begin{center}\rule{0.5\linewidth}{0.5pt}\end{center}

\section{按章节快速查找}\label{ux6309ux7ae0ux8282ux5febux901fux67e5ux627e}

{\def\LTcaptype{none} % do not increment counter
\begin{longtable}[]{@{}
  >{\raggedright\arraybackslash}p{(\linewidth - 2\tabcolsep) * \real{0.5000}}
  >{\raggedright\arraybackslash}p{(\linewidth - 2\tabcolsep) * \real{0.5000}}@{}}
\toprule\noalign{}
\begin{minipage}[b]{\linewidth}\raggedright
章节
\end{minipage} & \begin{minipage}[b]{\linewidth}\raggedright
直接查看该章参考
\end{minipage} \\
\midrule\noalign{}
\endhead
\bottomrule\noalign{}
\endlastfoot
第0章 前置基础 &
\href{../../chapters/00-prep/references.md}{references.md} \\
第1章 NumPy &
\href{../../chapters/01-numpy/references.md}{references.md} \\
第2章 SymPy &
\href{../../chapters/02-sympy/references.md}{references.md} \\
第3章 SciPy &
\href{../../chapters/03-scipy/references.md}{references.md} \\
第4章 Pandas &
\href{../../chapters/04-pandas/references.md}{references.md} \\
第5章 Matplotlib &
\href{../../chapters/05-matplotlib/references.md}{references.md} \\
第6章 NetworkX &
\href{../../chapters/06-networkx/references.md}{references.md} \\
第7章 Statsmodels &
\href{../../chapters/07-statsmodels/references.md}{references.md} \\
第8章 scikit-learn &
\href{../../chapters/08-sklearn/references.md}{references.md} \\
\end{longtable}
}

\backmatter
